% Introduction — DESI paper v3 (complete rewrite)
% Honest about what we test and what we conclude

\section*{Introduction}

The origin of anatomically modern \textit{Homo sapiens} during the Middle Pleistocene (~1.0--0.3 Ma) remains one of the most contested questions in human evolutionary biology. The Middle Pleistocene Transition (MPT; $\sim$1.25--0.70 Ma) is characterized by a sparse hominin fossil record yet represents the interval containing the last common ancestor of modern humans, Neanderthals, and Denisovans~\citep{Stringer2016}. Genomic data offer a principled window into this period, because the distribution of pairwise coalescence times across the genome encodes past population sizes, structure, and connectivity over millions of years.

The most precise recent inference of ancestral population dynamics during this interval came from FitCoal~\citep{Hu2023}, a composite-likelihood method applied to the site frequency spectrum (SFS) of 3,154 worldwide genomes. FitCoal inferred a severe panmictic bottleneck at approximately 930 ka, reducing ancestral human effective population size ($N_e$) to $\sim$1,280 individuals for approximately 117,000 years --- representing a reduction to roughly 1.3\% of the pre-bottleneck population. This interpretation was subsequently challenged on two fronts: \citet{Cousins2025b} showed that a panmictic model without a severe bottleneck fits the SFS comparably well, while \citet{Cousins2025} applied Cobraa --- a structured coalescent hidden Markov model --- and inferred two ancestral populations diverging at $\sim$1.5 Ma with $\sim$20\% gene flow at $\sim$300 ka, providing a structured alternative that does not require the catastrophic bottleneck. These competing analyses share a fundamental limitation: all are derived from the SFS, which cannot distinguish between a bottleneck and ancestral population structure, because both scenarios generate qualitatively similar rare-variant enrichment~\citep{Cousins2025b,Myers2008}.

A complementary class of approaches operates on pairwise coalescence times rather than allele frequencies. PSMC~\citep{Li2011} and MSMC2~\citep{Schiffels2014} infer $N_e(t)$ trajectories from single diploid genomes and pairs, respectively, and thus do capture deep genealogical signal; however, they estimate a single per-individual trajectory and do not provide a direct statistical test of specific demographic parameterizations across continental populations. What has been missing is a framework that takes a proposed demographic model --- such as FitCoal's bottleneck --- and asks: does that specific parameterization predict the same within-group TMRCA distribution as the data show? Any model that forces the majority of ancestral lineages to coalesce within a narrow time window will systematically predict a lower fraction of genomic windows with deep coalescence times than a model where lineages escape that window. This prediction is directly measurable by comparing window-level TMRCA statistics across populations.

Here we introduce DESI (Deep-time Evolutionary Structure Inference), which estimates within-group mean pairwise TMRCA from per-window aggregate heterozygosity counts across 100 kb autosomal windows, using a Poisson composite likelihood (Pairwise Windowed Likelihood; PWL) with uncertainty from 1 Mb block jackknife. Crucially, we use the same statistical framework to directly test FitCoal's specific bottleneck parameterization: we simulate 112 parameter combinations spanning bottleneck severity ($N_e^{\text{bn}} = 300$--$20{,}000$) and duration (30--300 ka), apply the identical per-window TMRCA analysis to each simulation, and compare the resulting summary statistics to the 1000 Genomes Project observations. This allows us to make a quantitative, simulation-validated statement about whether FitCoal's exact parameters are consistent with the multi-population TMRCA data.

Applying DESI to 240 phased whole genomes from five continental super-populations~\citep{Byrska2022}, we find that within-African mean TMRCA ($1{,}229.5 \pm 6.4$ ka) exceeds within-East-Asian ($858.8 \pm 5.4$ ka) by $370.7 \pm 8.4$ ka ($Z = 73$), a difference consistent across all 22 autosomes, unchanged in putatively neutral windows, and tracking known ancestry proportions in admixed populations. Our bottleneck parameter scan shows that FitCoal's exact parameterization ($N_e = 1{,}280$, 117 ka) predicts a fraction of windows with within-AFR TMRCA above 930 ka of 0.615, compared with the observed 0.707 ($Z = 4.6$, $p < 0.0001$), and underestimates mean within-AFR TMRCA by 10.2\%. The compatible parameter space requires a substantially weaker bottleneck ($N_e^{\text{bn}} \approx 2{,}000$), suggesting that if a bottleneck occurred, its severity has been overestimated by FitCoal.
